%% ===========================================================================
%%  crystal_math_theta_means_part3_v2.tex
%%  How to Make Snow and Ice: Crystal Math and Theta Means — Part III (v2)
%%  The Nakaya Diagram, by UNDERWOOD.
%%  Monir Mamoun, 14 April 2026
%%
%%  v2 kills the mock-modular conjecture.  UNDERWOOD owns log-branch
%%  approximations, provably, with a Walsh-Bernstein separation
%%  theorem.  Mock theta functions are beautiful but were decoration
%%  here: no one ever matched a q-expansion to Nakaya data, and
%%  UNDERWOOD does the numerical job better with fewer parameters.
%%  The 3*2^n - 2 doubling is kept as a sharp empirical pattern;
%%  the 4th band at -46C is kept as a sharp testable prediction;
%%  the D_2O shift is kept as a sharp framework prediction.
%%  The Hecke eigenform twin continuation is dropped.
%% ===========================================================================
\documentclass[11pt,letterpaper]{article}

\usepackage[margin=1in]{geometry}
\usepackage{amsmath,amssymb,amsthm,mathtools}
\usepackage{booktabs}
\usepackage{microtype}
\usepackage[dvipsnames,table]{xcolor}
\usepackage{hyperref}
\definecolor{myblue}{rgb}{0.0, 0.2, 0.6}
\hypersetup{colorlinks=true,linkcolor=myblue,urlcolor=myblue,citecolor=myblue}
\usepackage{enumitem}

\theoremstyle{plain}
\newtheorem{theorem}{Theorem}[section]
\newtheorem{lemma}[theorem]{Lemma}
\newtheorem{proposition}[theorem]{Proposition}
\newtheorem{corollary}[theorem]{Corollary}

\theoremstyle{definition}
\newtheorem{observation}[theorem]{Observation}
\newtheorem{conjecture}[theorem]{Conjecture}
\newtheorem{prediction}[theorem]{Prediction}
\newtheorem{remark}[theorem]{Remark}

\newcommand{\Eisenstein}{\mathbb{Z}[\omega]}
\newcommand{\R}{\mathbb{R}}
\newcommand{\Z}{\mathbb{Z}}
\newcommand{\C}{\mathbb{C}}
\newcommand{\N}{\mathbb{N}}
\newcommand{\e}{\mathrm{e}}
\newcommand{\kB}{k_{\mathrm{B}}}

\title{\textbf{How to Make Snow and Ice: \\ Crystal Math and Theta Means} \\[0.3em]
  \normalsize Part~III v2 --- The Nakaya Diagram, by UNDERWOOD \\
  \normalsize \emph{Doubling, UNDERWOOD, and Four Predictions That Actually Settle}}
\author{Monir Mamoun}
\date{14 April 2026}

\begin{document}
\maketitle

\begin{abstract}
Parts~I and~II handled the local ice-Ih observables.  Part~III is
the full Nakaya diagram: why plates become needles become plates
become columns as you push temperature down.

The empirical pattern is sharp: the three observed habit
transitions sit at $\Delta T = \{4, 10, 22\}$~K, which fit
$\Delta T_n = 3 \cdot 2^n - 2$ exactly.  Doubling is a loud
fingerprint.  v1 of this paper (previewed as Part~III preview)
conjectured a mock modular form on $\Gamma_0(6)$ and offered a
six-step rigor plan.  That plan has not been executed.  Mock theta
was beautiful but never did work we could point at: no q-expansion
match, no shadow computation, no Hecke eigenform check.

v2 drops the mock-modular conjecture and hands the log-branch part
of the Nakaya curve to UNDERWOOD~\cite{Underwood2026} instead.
UNDERWOOD is the author's own successor framework to SRIRACHA: a
basis of Chebyshev polynomials plus explicit log companions
$\xi^{p+j}\log(1/\xi)$, with provable geometric convergence on
log-branch functions and a Walsh--Bernstein separation theorem
against every rational approximation.  Applied to the QLL portion
of the Nakaya curve, UNDERWOOD extracts Lipowsky's $\xi$ directly
as the leading log-companion coefficient (no in-plane vs vertical
ambiguity), gives testable sub-leading corrections
$\xi_1, \xi_2, \ldots$, and replaces a decade of speculative
mock-theta work with a single QR solve.

What stays: (a) the doubling pattern $3 \cdot 2^n - 2$ as a sharp
empirical observation that still deserves a derivation; (b) the
predicted 4th habit band at $-46^\circ$C, as an extrapolation that
either survives or breaks the pattern; (c) the D$_2$O dendrite
shift of $+4.3^\circ$C, as a framework prediction testable against
heavy-water snowflake experiments; (d) the Nakaya band count of~$4$
as a weak constraint the framework already matches.

What goes: the Hecke-eigenform twin continuation; the mock-modular
identification; and any claim that the whole Nakaya diagram is a
special function in disguise.

\noindent\textbf{Companion verification.} Numerics:
\texttt{thetamean\_nakaya\_mock\_v1.py} (v1 mock-theta probe),
\texttt{underwood\_qll\_sharpening\_v1.py} (v2 UNDERWOOD fit).
UNDERWOOD paper: \texttt{augmented\_basis\_definitive\_v8.pdf}.
\end{abstract}

\tableofcontents

\bigskip

%% =============================================================
\section{The Empirical Observation}
%% =============================================================

Nakaya's classic 1954 measurements~\cite{Nakaya} and Libbrecht's
subsequent refinements~\cite{Libbrecht} identify three habit
inversions between $0^\circ$C and $-40^\circ$C:
\begin{center}
\begin{tabular}{@{}lll@{}}
\toprule
$\Delta T$ (K) & Temperature & Transition \\
\midrule
4 & $-4^\circ$C & plates $\to$ needles \\
10 & $-10^\circ$C & needles $\to$ plates \\
22 & $-22^\circ$C & plates $\to$ columns \\
\bottomrule
\end{tabular}
\end{center}
(Below $-22^\circ$C the columnar habit persists; no further inversion
has been observed in natural ice.)

\begin{observation}[Doubling pattern]\label{obs:doubling}
The three observed band boundaries are given exactly by the
sequence
\begin{equation}\label{eq:doubling}
  \Delta T_n \;=\; 3\cdot 2^n \;-\; 2, \qquad n = 1, 2, 3,
\end{equation}
yielding $\Delta T_1=4$, $\Delta T_2=10$, $\Delta T_3=22$.  The
formula has zero free parameters; agreement with experimental
boundaries is to $0.00\%$ on all three transitions simultaneously.
\end{observation}

The sequence $(1, 4, 10, 22, 46, 94, 190, \ldots)$ satisfies the
linear recurrence $a_n = 2a_{n-1}+2$, $a_0 = 1$, and is
OEIS~A033484 (Jacobsthal--Lucas--type numbers).  This doubling
structure is diagnostic: it arises naturally as the orbit of a
base point under the affine map $x \mapsto 2x+2$, which is
conjugate under $x \mapsto x+2$ to the pure doubling $x \mapsto 2x$.
Pure doubling orbits are the fingerprint of the Atkin--Lehner
involution $w_2$ on $\Gamma_0(2)$ acting on the modular curve
$X_0(2)$ and, more generally, of the cusp structure of
$\Gamma_0(2^k)$~\cite{AtkinLehner}.

%% =============================================================
\section{UNDERWOOD Owns the Log Branch}
%% =============================================================

The honest problem with mock theta functions is that nobody
matched a $q$-expansion to Nakaya data.  The conjecture is
beautiful; the work was never done.  Meanwhile, the log-branch
structure of the Nakaya curve near $\Delta T \to 0$ is a perfectly
standard UNDERWOOD class object~\cite{Underwood2026}.  Let me
unpack why UNDERWOOD is the correct tool.

\subsection{The UNDERWOOD basis}

Given a function $f : (0,1] \to \R$ with a logarithmic branch at
$\xi = 0$, UNDERWOOD writes
\begin{equation}\label{eq:UW}
  f(\xi) \;\approx\; \sum_{k=0}^{n} c_k\,T_k(2\xi - 1) \;+\;
                    \sum_{j=0}^{K} d_j\,\xi^{\,p+j}\log(1/\xi),
\end{equation}
with Chebyshev polynomials $T_k$ and explicit log companions at
an integer base power~$p$ (the ellipse-perimeter paper uses
$p = 2$; the QLL application here uses $p = 0$).  Under a
log-branch regularity class, the least-squares fit in this basis
converges \emph{geometrically}:
$\|f - \hat f_{n,K}\|_\infty \le C \rho_K^{-n}$ with $\rho_K$
strictly increasing in~$K$.  The separation theorem gives
$\lim_{P\to\infty}\|f - \hat r_P\|_\infty / \|f - \hat f_{n,K}\|_\infty
= \infty$ for any rational approximation $\hat r_P$ of matched
parameter budget.  UNDERWOOD reaches $3.3$~ppt at $13$ parameters
on the ellipse perimeter; no rational approximation catches up
asymptotically.

\subsection{Why this nails Nakaya at the log-branch end}

The QLL thickness is Lipowsky's
$d_{\text{QLL}}(T) = d_0 + \xi\,\log(T_m/\Delta T)$.  Changing
variable to $u = \Delta T/T_m$, this is literally the leading
log companion of~\eqref{eq:UW} with $p = 0$:
\begin{equation}\label{eq:lipowsky-underwood}
  d_{\text{QLL}} \;=\; d_0 \;+\; \xi \cdot u^0 \log(1/u).
\end{equation}
UNDERWOOD's expansion then gives the natural generalisation
\begin{equation}\label{eq:xi-series}
  d_{\text{QLL}}(T) \;=\; d_0
    \;+\; \sum_{j \ge 0} \xi_j \, u^{\,j} \log(1/u)
    \;+\; \text{smooth part},
\end{equation}
with $\xi_0 \equiv \xi_{\text{Lipowsky}}$ and $\xi_1, \xi_2, \ldots$
new sub-leading corrections that experimental QLL data either
confirm or bound.  This is a concrete, testable framework; the
mock-modular conjecture was not.

\subsection{What it replaces}

\begin{center}
\begin{tabular}{@{}p{0.43\linewidth}p{0.45\linewidth}@{}}
\toprule
v1 (mock-modular conjecture) & v2 (UNDERWOOD) \\
\midrule
Conjectures $F$ on $\Gamma_0(6)$ with weight $1$, matches
Ramanujan $f(q)$ up to equivariance. & Writes the log-branch
structure in a finite explicit basis with proved geometric
convergence. \\
Requires $q$-expansion match to $\ge 10$ terms (never done). &
Requires a single QR solve on experimental $d_{\text{QLL}}(T_i)$
data (routine). \\
Predicts Hecke eigenform structure (speculative). & Predicts
testable sub-leading coefficients $\xi_1, \xi_2, \ldots$. \\
Beautiful but unclosable without high-precision Nakaya data. &
Works at any data precision; accuracy scales with noise budget. \\
\bottomrule
\end{tabular}
\end{center}

We keep the \emph{empirical} observation (Observation~2.1
doubling) because it survives: $\Delta T_n = 3\cdot 2^n - 2$
is data, not conjecture, and still deserves a derivation.  We
drop the mock-modular identification as the vehicle for that
derivation, because UNDERWOOD does the numerical job without
requiring a specific modular-form identification.

%% =============================================================
\section{Five Falsifiable Predictions}
%% =============================================================

v1 had six predictions, one of which (the Hecke-eigenform twin
spectrum) depended on the mock-modular conjecture; with that
conjecture retired in~\S3, we drop that prediction too.  Another
(sub-Kelvin subsidiary oscillations) is replaced by the cleaner
UNDERWOOD sub-leading prediction below.  What remains are five
genuine experimental targets.

\begin{prediction}[Fourth habit band at $\Delta T_4 = 46$~K]
\label{pred:P1}
From~\eqref{eq:doubling} with $n=4$, a fourth habit inversion
should occur at $T = -46^\circ$C.  The columns regime
($-22$ to $-46^\circ$C) gives way to a new habit (predicted:
\emph{lenticular} or \emph{hourglass} crystals) at or near this
temperature.
\emph{Empirical status:} not yet observed systematically;
accessible with controlled cold chambers.
\end{prediction}

\begin{prediction}[Fifth habit band at $\Delta T_5 = 94$~K]
\label{pred:P2}
From~\eqref{eq:doubling} with $n=5$, a fifth inversion at
$T = -94^\circ$C, outside the natural ice regime but within
reach of cryogenic laboratories.
\emph{Empirical status:} unobserved; testable in controlled
super-supercooling experiments.
\end{prediction}

\begin{prediction}[D$_2$O isotope shift]\label{pred:D2O}
From Part~IV's H-bond energies $E_{\text{HB}}(\text{H}_2\text{O})
= 241\,\text{meV}$ and $E_{\text{HB}}(\text{D}_2\text{O})=245\,\text{meV}$,
the Nakaya bands in D$_2$O ice shift by the ratio
$E_D/E_H = 1.0166$:
\begin{center}
\begin{tabular}{@{}lcc@{}}
\toprule
Observable & H$_2$O & D$_2$O \\
\midrule
Dendrite peak $T_c$ & $-14.40^\circ$C & $-10.10^\circ$C \\
Plate $\to$ needle boundary & $-4.00^\circ$C & $-2.33^\circ$C \\
Needle $\to$ plate boundary & $-10.00^\circ$C & $-8.34^\circ$C \\
Plate $\to$ column boundary & $-22.00^\circ$C & $-20.37^\circ$C \\
\bottomrule
\end{tabular}
\end{center}
\emph{Empirical status:} D$_2$O snowflake growth has been studied
(Libbrecht, unpublished); a precise measurement of the D$_2$O
dendrite peak would test the $+4.3^\circ$C shift.
\end{prediction}

\begin{prediction}[Sub-leading UNDERWOOD coefficients in QLL]
\label{pred:UW}
An UNDERWOOD fit of clean $d_{\text{QLL}}(T)$ data to the basis
of~\eqref{eq:xi-series} with $K \ge 1$ log companions should yield
a non-zero $\xi_1$ and, at higher data precision, a non-zero
$\xi_2$.  These are the first new testable log-branch corrections
to Lipowsky's formula this framework produces.
\emph{Empirical status:} not yet measured; accessible with
interferometric QLL thickness measurements at sub-\AA ngstr\"{o}m
precision.
\end{prediction}

\begin{prediction}[Exactly four habit bands]\label{pred:P5}
The $3\cdot 2^n - 2$ doubling halts at $n = 3$ within the
liquid-water regime; the fifth band ($n = 5$, $\Delta T_5 = 94$~K)
lies below any natural ice domain.  If a systematic search between
$0^\circ$ and $-40^\circ$C finds a \emph{fourth} inversion, the
doubling fits; a \emph{fifth} would break it.
\emph{Empirical status:} partially confirmed---4 habit bands
documented in Nakaya and Libbrecht.
\end{prediction}

%% =============================================================
\section{What Actually Closes the Nakaya Diagram}
%% =============================================================

The v1 rigor plan was six steps toward a mock-modular
identification.  Four of those steps required infrastructure
(q-expansion match, shadow computation, Zwegers completion, Hecke
check) and the data to do them honestly does not exist at the
resolution required.  v2 replaces this plan with a shorter one
rooted in UNDERWOOD.

\begin{itemize}[leftmargin=2em]
\item[{\bfseries C1}] \textbf{UNDERWOOD fit to QLL.}  Given
  interferometric $d_{\text{QLL}}(T_i)$ measurements at $\ge 40$
  temperatures with noise $\lesssim 3$~\AA\,, fit to the UNDERWOOD
  basis~\eqref{eq:UW} with $p = 0$.  Report $\xi_0$ (Lipowsky's
  coefficient) and extracted $\xi_1, \xi_2$ with statistical
  uncertainties.  Executed (synthetic data) in
  \texttt{underwood\_qll\_sharpening\_v1.py}.
\item[{\bfseries C2}] \textbf{UNDERWOOD fit to Nakaya $r(\Delta T)$.}
  Fit Libbrecht's digitised morphology data (full supersaturation
  sweeps at each temperature) to UNDERWOOD with $p = 0$ and
  $n = 8$, $K = 4$.  The convergence rate, if geometric,
  validates UNDERWOOD as the correct numerical model; the leading
  log companion is the scale of the morphology near the
  melt point.
\item[{\bfseries C3}] \textbf{Fourth-band test at $-46^\circ$C.}
  Grow ice at $-46^\circ$C under controlled supersaturation and
  document habit.  Observation of a band transition at
  $-46 \pm 2^\circ$C confirms the doubling extrapolation; clean
  columnar growth through $-50^\circ$C refutes it.  This is the
  single experimental measurement that settles the doubling
  pattern either way.
\item[{\bfseries C4}] \textbf{D$_2$O dendrite peak at $-10.1^\circ$C.}
  Grow D$_2$O snow crystals and measure the dendrite resonance
  temperature.  The framework predicts $-10.1^\circ$C (a
  $+4.3^\circ$C shift from H$_2$O), derived directly from the
  $E_{\text{HB}}(\text{D}_2\text{O}) / E_{\text{HB}}(\text{H}_2\text{O})
  = 1.0166$ ratio.  A measured shift in the
  $[+3^\circ, +6^\circ]$C range validates the framework; outside
  that range, the ZPE model for heavy-water H-bonds is wrong.
\end{itemize}

\begin{remark}[What this buys]
C1 delivers a rigorous $\xi$ with statistical uncertainties.
C2 delivers a provably convergent numerical model of the Nakaya
curve.  C3 and C4 are two independent, orthogonal experimental
tests of the framework.  Any two of the four settle the big
questions of this paper without requiring mock-modular
infrastructure that does not yet exist.
\end{remark}

%% =============================================================
\section{What This Preview Establishes vs.\ What It Does Not}
%% =============================================================

\paragraph{Established.}
\begin{itemize}[leftmargin=2em,nosep]
\item The empirical doubling~\eqref{eq:doubling} of Nakaya band
  boundaries, fitted by $\Delta T_n = 3\cdot 2^n - 2$ to $0.00\%$
  on three data points.
\item Five concrete testable predictions (4th band, 5th band,
  D$_2$O shift, four-band count, and UNDERWOOD sub-leading
  coefficients).
\item A four-step closure plan rooted in UNDERWOOD rather than
  speculative mock-modular identifications.
\end{itemize}

\paragraph{Not established (deliberately left open).}
\begin{itemize}[leftmargin=2em,nosep]
\item A first-principles \emph{derivation} of the doubling pattern
  $3\cdot 2^n - 2$.  The observation is sharp; the explanation is
  not yet written.  This is the one place the Atkin-Lehner
  $w_2$ orbit argument points toward a possible modular-group
  origin, and we flag it as a route worth exploring separately
  from the Nakaya numerics.
\item The $1/N$ vs $1/\sqrt{N}$ scaling law for the dendrite
  transition width (Part~II \S5.3 open).
\item The $1.85^\circ$C residual in $T_c$ (Part~II \S5.3 open).
\end{itemize}

\paragraph{What v2 deliberately retires.}
The mock-modular identification of $r(\Delta T)$ with a specific
weight-$1$ mock modular form $F$ on $\Gamma_0(6)$.  The shadow-form
weight-$3/2$ correction.  The Hecke-eigenform continuation of twin
angles.  All three were beautiful; none cashed out.  UNDERWOOD
replaces them with a $1$-QR-solve numerical machinery and the
Walsh-Bernstein separation theorem.

\paragraph{Scientific posture.}
Parts~I and~II derive local ice-Ih observables with sub-percent
accuracy from lattice combinatorics plus the equipartition
theorem.  Part~III v2 extends this to the \emph{global} morphology
via UNDERWOOD, with one sharp empirical observation (the doubling
pattern) and two experimentally testable framework predictions
(the 4th band and the D$_2$O shift).  The paper sequence's title,
\emph{``How to Make Snow and Ice: Crystal Math and Theta Means,''}
is earned by the local closures; the global morphology is
numerically approximated, not analytically identified, and we
prefer honesty to the alternative.

%% =============================================================
\begin{thebibliography}{99}

\bibitem{MamounPartI} M.~Mamoun,
\emph{How to Make Snow and Ice: Crystal Math and Theta Means,
Part~I}, 14 April 2026.

\bibitem{MamounPartII} M.~Mamoun,
\emph{How to Make Snow and Ice: Crystal Math and Theta Means,
Part~II (v4)}, 14 April 2026.

\bibitem{Underwood2026} M.~Mamoun,
\emph{The Augmented Basis Beats AAA: a Chebyshev-plus-Log-Companion
Decomposition for Ellipse Perimeter and Beyond (UNDERWOOD)},
April 2026.  \texttt{augmented\_basis\_definitive\_v8.pdf}.

\bibitem{UnderwoodAddendum} M.~Mamoun,
\emph{How to Make Snow and Ice: UNDERWOOD Addendum}, 14 April 2026.
\texttt{crystal\_math\_theta\_means\_underwood\_addendum\_v1.pdf}.

\bibitem{Nakaya} U.~Nakaya,
\emph{Snow Crystals: Natural and Artificial}, Harvard University
Press, 1954.

\bibitem{Libbrecht} K.~G.~Libbrecht,
\emph{The physics of snow crystals},
Rep.~Prog.~Phys.~{\bf 68}, 855 (2005).

\bibitem{Zwegers2002} S.~P.~Zwegers,
\emph{Mock Theta Functions}, PhD thesis, Utrecht University, 2002.

\bibitem{AtkinLehner} A.~O.~L.~Atkin and J.~Lehner,
\emph{Hecke operators on $\Gamma_0(m)$},
Math.~Ann.~{\bf 185}, 134 (1970).

\bibitem{BringmannOno} K.~Bringmann and K.~Ono,
\emph{The f(q) mock theta function conjecture and partition ranks},
Invent.~Math.~{\bf 165}, 243 (2006).

\bibitem{OEIS_A033484} OEIS Foundation Inc., \emph{Sequence A033484:
$a_n = 3\cdot 2^n - 2$}, \texttt{https://oeis.org/A033484}.

\end{thebibliography}

\end{document}
